\subsection*{Partie A}

\begin{enumerate}
	\item La variable aléatoire $X$ compte le nombre de succès d'une répétition de 10 épreuves de Bernoulli. Les épreuves sont indépendantes deux à deux. La probabilité du succès est $p=\np{0.4}$. Donc $X$ suit une loi binomiale de paramètres 10 ; \np{0.4} : $X \sim \mathcal{B}(10~;~\np{0.4})$.
	\item
		\begin{enumerate}
			\item $P(X = 2) = \binom{10}{2}\times p^2\times(1-p)^{10-2}= \binom{10}{2}\times\np{0.4}^{2}\times \np{0.6}^{8} \approx \np{0.1209}$
			\item Comme les évènements $\{X = i\}$ et $\{X = j\}$ sont disjoints pour $i \neq j$ :
\[
\begin{aligned}
P(X \leqslant 2)
&= P(\{X = 0\} \cup \{X = 1\} \cup \{X = 2\}) \\
&= \sum_{i=0}^{2} P(X=i) \\
&= \sum_{i=0}^{2}\binom{10}{i} p^{i}(1-p)^{10-i} \\
&\approx 0{,}1673
\end{aligned}
\]
			
			La probabilité que le phénomène El Niño soit dominant au plus deux années sur une période de 10~ans est d'environ \np{0.1673}.
		\end{enumerate}
	\item Comme $X$ suit une loi binomiale $\mathcal{B}(10~;~\np{0.4})$ : $E(X) = n\times p = 10 \times \np{0.4} = 4$. En moyenne le phénomène El Nino est dominant quatre années sur une période de 10~ans.
\end{enumerate}

\subsection*{Partie B}

\begin{enumerate}
	\item  \ {}
	
\begin{center}
\pstree[treemode=R,levelsep=3cm]{\TR{}}
{\pstree{\TR{$E_n~$}\taput{$p_n$}}
	{\TR{$E_{n+1}$} \taput{0,5}
	\TR{$\overline{E_{n+1}}$} \tbput{0,5}
	}
\pstree{\TR{$\overline{E_n}~$}\tbput{$1 - p_n$}}
	{\TR{$E_{n+1}$} \taput{0,3}
	\TR{$\overline{E_{n+1}}$} \tbput{0,7}
	}
}
\end{center}

	\item $p_0 = 0$\\ $1-p_0 = 1$
	
	$p_1 = P\left(E_1\right) = P\left(E_0\right)\times P_{E_{0}}\left(E_1\right) + P\left(\overline{E_0}\right)\times P_{\overline{E_{0}}}\left(E_1\right)
	= 0 \times \np{0.5} + 1 \times  \np{0.3} =   \np{0.3}$
	\item $\begin{gathered}p_{n+1} = P\left(E_{n+1}\right) = P\left(E_n\right)\times P_{E_{n}}\left(E_{n+1}\right) + P\left(\overline{E_n}\right)\times P_{\overline{E_{n}}}\left(E_{n+1}\right)\\
	= p_n \times \np{0.5} + (1-p_n) \times  \np{0.3} = \np{0.2}\times p_n + \np{0.3}\end{gathered}$
	\item
		\begin{enumerate}
			\item En utilisant la calculatrice la suite $\left(p_n\right)$ semble être croissante et avoir pour limite $\np{0.375}$.


\begin{center}
\begin{tabular}{cl}
\toprule
$n$ & \multicolumn{1}{c}{$p_n$} \\
\midrule
0 &\np{0} \\
1 &\np{0.3} \\
2 &\np{0.36} \\
3 &\np{0.372} \\
\vdots & \multicolumn{1}{c}{\vdots} \\
13 &\np{0.3749999996928} \\
14 &\np{0.37499999993856} \\
\bottomrule
\end{tabular}
\end{center}

			
			\item On formule l'hypothèse de récurrence $\mathcal{H}(n)~:~p_n\leqslant \np{0.375}$.
			
\textbf{On vérifie que $\mathcal{H}(0)$ est vraie (initialisation).}

Comme $p_0 = 0 \leqslant \np{0.375}$,  $\mathcal{H}(0)$ est vraie.

\textbf{Montrons que $\mathcal{H}(n)\Rightarrow \mathcal{H}(n+1)$ (hérédité).}

On suppose que $\mathcal{H}(n)$ est vraie, avec $n \in \N$.

$p_{n+1} = \np{0.2}\times p_n + \np{0.3}$

D'après $\mathcal{H}(n)$, $p_n \leqslant \np{0.375}$, en multipliant par $\np{0.2}$ on obtient :

$ \np{0.2}\times p_n \leqslant \np{0.075}$, en ajoutant  $\np{0.3}$ on obtient :

$ \np{0.2}\times p_n + \np{0.3} \leqslant \np{0.375}$, ainsi on obtient :

$p_{n+1} \leqslant \np{0.375} $ ($\mathcal{H}(n+1)$ est vérifiée)

\textbf{D'après le principe de récurrence},

$\forall n\in \N, n \geqslant 0, p_n\leqslant \np{0.375}$
		\item Pour tout $n\in \N$, $p_n$ désignant une probabilité est un nombre positif.

Soit ${n\in\N}$

$p_{n+1}-p_{n} = \np{-0.8}\times p_n + \np{0.3} $

$p_{n+1}-p_{n} \geqslant \np{-0.8}\times \np{0.375} + \np{0.3} $

$p_{n+1}-p_{n} \geqslant 0 $

Ceci étant valable quelque soit le choix de $n$, $\forall n \in \N,~p_{n+1}-p_{n} \geqslant 0 $, donc $\left(p_n\right)$ est croissante.
			\item $\left(p_n\right)$ est croissante et majorée donc $\left(p_n\right)$ est convergente vers $l$. De plus ${l\leqslant \np{0.375}}$.
		\end{enumerate}
	\item
		\begin{enumerate}
			\item \ {}
			
			\renewcommand{\arraystretch}{2}
			$\begin{array}{rcl}
			u_{n+1} 	&=& p_{n+1} - \dfrac{3}{8}\\
					&=& \np{0.2}\times p_n + \np{0.3} - \dfrac{3}{8}\\
					&=& \np{0.2}\times p_n +\np{0.2}\times \dfrac{3}{2} - \np{0.2}\times\dfrac{15}{8}\\
					&=& \np{0.2}\times p_n +\np{0.2}\times\left( \dfrac{3}{2} - \dfrac{15}{8}\right)\\
					&=& \np{0.2}\times p_n +\np{0.2}\times\left(  - \dfrac{3}{8}\right)\\
					&=& \np{0.2}\times \left(p_n - \dfrac{3}{8}\right)\\
			\end{array}$

Finalement quel que soit $n$ naturel, $u_{n+1} = \np{0.2}\times u_{n}$ : ceci prouve que $\left(u_n\right)$ est une suite géométrique de raison $\np{0.2}$.

			$u_0 = p_{0} - \dfrac{3}{8} = 0 - \dfrac{3}{8} = -\dfrac{3}{8}$

			Le premier terme est donné par : ${u_0 = - \dfrac{3}{8}}$

			\item Comme $\left(u_n\right)$ est une suite géométrique de raison $\np{0.2}$ de premier terme ${u_0 = - \dfrac{3}{8}}$ :

			$\forall n \in \N~:~u_n = - \dfrac{3}{8} \times \np{0.2}^{n} $

			Comme $\forall n \in \N~:~u_n = p_{n} - \dfrac{3}{8}$,

			$\forall n \in \N~:~p_n = -u_{n} + \dfrac{3}{8} = \left(-\np{0.2}^{n} +1 \right)\times \dfrac{3}{8} = \dfrac{3}{8}\left(1-\np{0.2}^{n}  \right)$
			
			\item
			Comme $\np{0.2}<1$, $\displaystyle\lim_{n\rightarrow +\infty}\np{0.2}^{n} = 0$ et donc $\displaystyle\lim_{n\rightarrow +\infty}p_{n} = \dfrac{3}{8}$
			\item La probabilité d'observer un phénomène El Niño dominant tend vers $\dfrac{3}{8}$ quand le nombre d'années d'observation augmente.
		\end{enumerate}	
\end{enumerate}
